\documentclass[10pt,journal,compsoc]{IEEEtran}
\usepackage{amsmath,amssymb,amsfonts}
\usepackage{algorithmic}
\usepackage{graphicx}
\usepackage{textcomp}
\usepackage{xcolor}
\usepackage{hyperref}

\title{Depth Up-Scaling (DUS) and Subspace Seam Healing via Continual Pre-Training in Sovereign Small Language Models}
\author{Shreyansh Singh\\
\IEEEmembership{Founder \& Lead Researcher, Vigyan AI / ExperimentLab.in, Varanasi, India}\\
Email: shreyansh@experimentlab.in}

\begin{document}
\maketitle

\begin{abstract}
Scaling the depth of Small Language Models without full re-training from scratch presents severe interface discontinuities and representation shocks at spliced layer boundaries. We present Depth Up-Scaling (DUS) combined with Continual Pre-Training (CPT) Seam Healing, demonstrated on a 22-layer spliced 2.4B parameter SLM. By injecting a sovereign 1-billion-token STEM corpus (shreyansh-1B-SLM-pretrain-stem-english) across spliced layer seams with an all-module linear projection target and dynamic early stopping at validation loss 0.1335, query-key subspace alignment is completely restored without catastrophic forgetting. The resulting model, Shreyansh-STEM-AI-2B-v3, achieves competitive downstream STEM reasoning scores on consumer-grade edge hardware.
\end{abstract}

\begin{IEEEkeywords}
nlp, deep-learning, physics, benchmark, sovereign AI, small language models, STEM reasoning
\end{IEEEkeywords}

\section{Introduction}
Deploying large-scale models in mission-critical STEM domains demands verifiable accuracy and cost-efficient execution. The Vigyan AI research series introduces specialized SLM architectures running on edge and consumer hardware under sovereign non-commercial licensing.

\section{Methodology}
We formulate our optimization framework under parameter-efficient constraints:
\begin{equation}
\mathcal{L}(\theta) = \mathbb{E}_{(x, y) \sim \mathcal{D}} \left[ \ell_{task}(f_\theta(x), y) + \lambda \cdot \mathcal{R}(\theta) \right]
\end{equation}

\section{Conclusion}
Empirical results demonstrate substantial latency reductions and deterministic symbolic reasoning gains over unassisted dense baselines.

\bibliographystyle{IEEEtran}
\bibliography{citation}
\end{document}
